\documentclass{standalone}
\usepackage{circuitikz}
\begin{document}
\begin{circuitikz}[american, scale=1.0, transform shape]
    % Linee di alimentazione
    \draw (-0.5, 5) node[left]{$V_{CC}$} to[short, o-] (8.5, 5);
    \draw (-0.5, 0) node[left]{GND} to[short, o-] (8.5, 0);

    % Ingresso
    \draw (-0.5, 2.5) node[left]{$v_{in}$} to[C, l=$C_{in}$, o-*] (1.5, 2.5);
    
    % Rete di polarizzazione del CE
    \draw (1.5, 5) to[R, l=$R_1$, *-] (1.5, 2.5) to[R, l=$R_2$, -*] (1.5, 0);
    
    % Transistor Q1 (CE)
    \draw (1.5, 2.5) to[short] (2.5, 2.5) node[npn, anchor=B](Q1){$Q_1$};
    
    % Connessioni di collettore ed emettitore per CE
    \draw (Q1.C) to[R, l=$R_C$, *-*] (Q1.C |- 0,5);
    \draw (Q1.E) to[R, l=$R_{E1}$, -*] (Q1.E |- 0,0);
    
    % Accoppiamento diretto (DC) al CC
    \draw (Q1.C) to[short] ++(2,0) node[npn, anchor=B](Q2){$Q_2$};
    
    % Connessioni CC
    \draw (Q2.C) to[short, -*] (Q2.C |- 0,5);
    \draw (Q2.E) to[R, l=$R_{E2}$, *-*] (Q2.E |- 0,0);
    
    % Uscita
    \draw (Q2.E) to[short] ++(1,0) to[C, l=$C_{out}$, -o] ++(1.5,0) node[right]{$v_{out}$};
    
    % Riquadri per evidenziare gli stadi
    \draw[dashed, blue!50, rounded corners, thick] (0.5, -0.5) rectangle (3.5, 5.5);
    \node[font=\small\bfseries, blue] at (2, 5.8) {Stadio CE ($|A_V| \gg 1$)};
    
    \draw[dashed, red!50, rounded corners, thick] (3.8, -0.5) rectangle (6.5, 5.5);
    \node[font=\small\bfseries, red] at (5.15, 5.8) {Stadio CC (Bassa $Z_{out}$)};
\end{circuitikz}
\end{document}
